\documentclass[11pt]{amsart}
\usepackage[margin=1.15in]{geometry}
\usepackage{amsmath,amssymb,amsthm,mathtools,hyperref}
\newtheorem{theorem}{Theorem}
\newtheorem{lemma}[theorem]{Lemma}
\newtheorem{proposition}[theorem]{Proposition}
\newtheorem{corollary}[theorem]{Corollary}
\theoremstyle{definition}\newtheorem{definition}[theorem]{Definition}
\newcommand{\Fp}{\mathbb F_p}
\newcommand{\C}{\mathbb C}
\renewcommand{\Vec}{\operatorname{Vec}}
\newcommand{\Rep}{\operatorname{Rep}}
\newcommand{\FPdim}{\operatorname{FPdim}}
\newcommand{\Inv}{\operatorname{Inv}}
\newcommand{\PW}{\operatorname{PW}}
\newcommand{\cB}{\mathcal B}
\newcommand{\cC}{\mathcal C}
\newcommand{\cD}{\mathcal D}
\newcommand{\cM}{\mathcal M}
\newcommand{\cA}{\mathcal A}
\title{Modular categories of dimension $p^6$ at odd primes}
\author{Antares}
\date{14 November 2026}
\begin{document}
\begin{abstract}
For every odd prime $p$, we classify the nonpointed complex modular categories
of Frobenius--Perron dimension $p^6$ and multiplicative central charge $-1$.
There are exactly $4p+4$ braided equivalence classes: $2p+4$ products of a prime
category of dimension $p^5$ with a pointed category of dimension $p$, and $2p$
prime categories. We give explicit cocycles and central embeddings realizing
all classes, compute the complete equivalence action, and determine their
invertible groups and Frobenius--Schur exponents. Together with the known
charge-one and pointed classifications, this completes the classification
in dimension $p^6$.
\end{abstract}
\maketitle

\section{Scope and classification}
Throughout $p$ is an odd prime, $\epsilon\in\Fp^\times$ is a fixed nonsquare,
$\zeta=\exp(2\pi i/p)$, and $\eta=\exp(2\pi i/p^2)$.
All categories are complex fusion categories. A modular category includes a
ribbon structure. Write $\xi$ for its multiplicative central charge.
Categories of odd prime-power Frobenius--Perron dimension are integral and
pseudo-unitary; see \cite[Propositions 8.23--8.24 and Theorem 8.28]{ENO05}.
Their ribbon structure is unique: the ambiguity of ribbon structures in a nondegenerate
braided category is its group of invertibles of order two, which here is trivial.
Thus the given ribbon structure is the positive one and braided equivalences
preserve $\xi$.

\begin{theorem}\label{thm:main}
There are exactly $4p+4$ braided equivalence classes of nonpointed modular
categories $\cC$ with $\FPdim\cC=p^6$ and $\xi(\cC)=-1$.
They consist of the following pairwise distinct classes.
\begin{enumerate}
\item Product classes $\cC_0(c,Q)$, where $c\in\{1,\epsilon\}$ and
\[
 Q\in\{0,\ X^2,\ \epsilon X^2,\ X^2+dY^2\ (d\in\Fp^\times)\}.
\]
Each is the product of a prime modular category of dimension $p^5$ and a
pointed modular category of dimension $p$.
\item Prime classes $\cC_1(c,s)$, where $c\in\{1,\epsilon\}$ and $s\in\Fp$.
Their invertible group is $C_{p^2}\times C_p$; every invertible of order $p$
has twist one.
\end{enumerate}
All nonpointed categories in this theorem have $p^3$ invertible simples.
The constructions and complete equivalence criteria are given below.
\end{theorem}

The prior classification machinery is substantial. DGNO's anisotropic-core
and maximal-isotropic theory \cite{DGNO}, ENO's extension obstruction theory
\cite{ENO}, and the associator-equivalence correction of Davydov--Nikshych
\cite{DN} are inputs. Crossed pointed categories and their comparison cochains
were classified by Naidu \cite{Naidu}; compatible central lifts and cleft
quotients were studied by Mason--Ng \cite{MN}, and the pointed-centre Morita
criterion by Naidu--Nikshych \cite{NN}. We use the centralizer model of
Gannon--Riesen \cite{GR}. The new content is the explicit all-prime normal forms,
their complete unmarked equivalence law, and the prime/product consequences.

Two close classifications have different scopes. Green--Nikshych
\cite[Section 7.1, Proposition 7.6]{GN} classify the $p+1$ integral nonpointed
\emph{reductive} categories of dimension $p^6$; these are metabolic centres.
The categories of Theorem~\ref{thm:main} instead have a unique nontrivial maximal
Tannakian subcategory. Schopieray \cite[Section 6]{Schopieray} gives the
extraspecial-centre realization underlying the dimension-$p^5$ setting, while
leaving its odd-prime braided classification open there. Neither result gives
the parameter orbits or prime dimension-$p^6$ family below.

\section{Prior reductions and a canonical subcategory}
The following are prior results. The maximal-isotropic theory of DGNO
\cite{DGNO} identifies the anisotropic metric core of an integral nilpotent
modular category. If $\FPdim\cC=p^6$ and $\xi(\cC)=-1$, a maximal Tannakian
subcategory $E$ has dimension $p^2$, and its local de-equivariantization is an
anisotropic pointed modular category $\cA=\cC(A,q)$ of dimension $p^2$.
Here necessarily $A=C_p^2$. Its full de-equivariantization is a braided
$G$-crossed category, where $E=\Rep(G)$ and $|G|=p^2$.
The orthogonal group of an anisotropic plane has order $2(p+1)$, prime to $p$,
so the action $G\to O(A,q)$ is trivial.
Gannon--Riesen \cite[Theorem 4.1]{GR} then makes the full de-equivariantization
pointed, say $\Vec_\Gamma^\omega$, with a central extension
\[
1\longrightarrow A\longrightarrow\Gamma\longrightarrow G\longrightarrow1.
\]
Its centre has the specified factorization
\begin{equation}\label{eq:stabilize}
Z(\Vec_\Gamma^\omega)\simeq\cA^{\rm rev}\boxtimes\cC.
\end{equation}
The complement is $\cA^{\rm rev}$, not $\cA$; this sign convention will matter
in the explicit formulas.
Czenky--Plavnik--Schopieray \cite[Lemma 2]{CPS} give
\begin{equation}\label{eq:invbound}
p^2\leq |\Inv\cC|\leq p^3
\end{equation}
for a nonpointed modular category of dimension $p^6$.

\begin{lemma}\label{lem:flux}
If an anisotropic pointed braided category $\cC(B,q_B)$ embeds braidedly in
$Z(\Vec_K^\omega)$, its forgetful degree map $B\to K$ is injective and has
central image.
\end{lemma}
\begin{proof}
An invertible central object has an invertible underlying degree $k_b$;
half-braiding with every degree forces $k_b\in Z(K)$. Degrees multiply under
tensor product. A kernel object lies over the unit and its self-braiding is one,
by the unit axiom and naturality of a half-braiding. Thus $q_B(b)=1$, and
anisotropy forces $b=0$.
\end{proof}
Consequently every pointed presentation of \eqref{eq:stabilize} has degree
group of order $p^4$ with a central injected $C_p^2$, hence has class at most two.

\begin{lemma}\label{lem:canonical}
If $\cC$ is nonpointed with dimension $p^6$ and charge $-1$, then $G=C_p^2$,
$|\Inv\cC|=p^3$, and $E$ is the unique maximal Tannakian subcategory.
Its simples are exactly the invertibles of twist one.
\end{lemma}
\begin{proof}
If $G=C_{p^2}$, its central extension $\Gamma$ is abelian. The crossed object
action is conjugation by a lift, hence fixes every object. All complex projective
irreducibles of cyclic $G$ are one-dimensional, so the equivariantization would
be pointed. Thus $G=C_p^2$.
For each fixed neutral simple $a\in A$ let $\alpha_a\in H^2(G,\C^\times)$ be
its equivariance obstruction. Tensor coherence gives
$[\alpha_{a+b}]=[\alpha_a][\alpha_b]$. Since $H^2(C_p^2,\C^\times)=C_p$,
the kernel has order at least $p$. Each kernel element has $p^2$ invertible
equivariant lifts. The bound \eqref{eq:invbound} forces equality, a kernel line
$M\subset A$, and exhaustion of the invertibles by these lifts.
Their twists are the pullbacks of $q|_M$. Anisotropy makes the twist-one
invertibles exactly the $p^2$ lifts of the neutral unit, namely $E$.
Every maximal Tannakian subcategory has dimension $p^2$ by the anisotropic-core
reduction, and is pointed since every group of order $p^2$ is abelian. It must
therefore equal $E$.
\end{proof}
Thus braided equivalences preserve $E$ and induce equivalences of the full
crossed categories, allowing an automorphism of $G$ and a metric automorphism
of $A$. Generic pointed Morita equivalence is not required for the nonpointed
branch after this reduction.

\section{The obstruction and equivalences of its trivializations}
Write
\[
H^*(G,\Fp)=\Lambda(x_1,x_2)\otimes\Fp[y_1,y_2],\qquad y_i=\beta(x_i),
\]
where $\beta$ is the mod-$p$ Bockstein. Put $e=x_1x_2$. A trivial-action class
has unique coordinates
\begin{equation}\label{eq:kappa}
\kappa=C e+P_1y_1+P_2y_2\in H^2(G,A).
\end{equation}
Here $C$ is the commutator vector and $P_i$ are the power vectors of $\Gamma$.
Write $q(a)=\zeta^{Q_A(a)}$ and
$B_A(a,b)=Q_A(a+b)-Q_A(a)-Q_A(b)$.

\begin{proposition}\label{prop:obstruction}
The associativity obstruction vanishes if and only if
$B_A(C,P_1)=B_A(C,P_2)=0$.
For a nonpointed equivariantization $C\ne0$. Thus its power map has rank zero
or one, with image contained in the line $C^\perp\ne\langle C\rangle$.
\end{proposition}
\begin{proof}
ENO's odd-order Pontryagin--Whitehead formula
\cite[Proposition 8.13]{ENO} is the image of $Q_A(\kappa)$ under
$\Fp\to\C^\times$, $1\mapsto\zeta$.
The three polynomial classes $y_1^2,y_1y_2,y_2^2$ map to zero in
$H^4(G,\C^\times)$. The two remaining basis elements $ey_i$ have independent
Bocksteins
\[
\beta(ey_1)=y_1^2x_2-x_1y_1y_2,\qquad
\beta(ey_2)=y_1y_2x_2-x_1y_2^2.
\]
These identify the two-dimensional $H^4(G,\C^\times)$, and their coefficients
in $Q_A(\kappa)$ are $B_A(C,P_i)$. This proves the criterion.
If $C=0$, then $\Gamma$ is abelian and the action fixes all $p^4$ simples.
Their obstruction homomorphism to $H^2(G,\C^\times)=C_p$ has kernel of order
at least $p^3$, producing at least $p^5$ invertibles, contrary to
\eqref{eq:invbound}. Finally anisotropy gives the stated orthogonal line.
\end{proof}

\subsection{The $H^1$ identification}
It is essential to include the correction in
\cite[Corollaries 2.25, 2.27 and Remark 8.8]{DN}:
associator equivalence classes form a torsor over a cokernel, not automatically
over all of $H^3(G,\C^\times)$.

\begin{proposition}\label{prop:h1}
For the trivial object action on the odd metric group $A$, the automorphism
$\lambda\in H^1(G,A)$ of the lower extension data acts on the associator
torsor by translation by
\begin{equation}\label{eq:h1}
B_A(\kappa,\lambda)\in H^3(G,\C^\times).
\end{equation}
Here the pairing is applied to cup product, followed by $1\mapsto\zeta$.
\end{proposition}
\begin{proof}[Proof of the full action]
The identity-component Postnikov stage of $B\operatorname{Pic}(\cA)$ has
$\pi_2=A$, $\pi_3=\C^\times$ and quadratic invariant $q$
\cite[Propositions 7.3 and 7.6]{ENO}; equivalently it is
the homotopy fibre of the quadratic map
$K(A,2)\to K(\C^\times,4)$. This is the odd-order quadratic model underlying
ENO's formula. A loop of the lower map, represented by $\lambda$, gives the
class $\kappa+u\lambda$ on $BG\times S^1$, with $u$ the generator of
$H^1(S^1,\mathbb Z)$. Polarization and $u^2=0$ give
\[
q(\kappa+u\lambda)=q(\kappa)+uB_A(\kappa,\lambda).
\]
Slanting with the circle gives \eqref{eq:h1}, precisely the comparison cube
of \cite[Proposition 2.24]{DN}. This computation must be retained when
passing from fixed products to equivalences of crossed extensions.
\end{proof}

The group $H^3(G,\C^\times)$ is the space of binary quadratic forms. In the
coordinates above its basis is represented by $x_1y_1,x_1y_2,x_2y_2$; their
Bocksteins are $y_1^2,y_1y_2,y_2^2$.
If $P_i=0$, \eqref{eq:h1} is zero, since $e\lambda=0$ in cohomology.
If $\kappa=C e+T y_1$, with $Q_A(T)\ne0$, its image is the two-dimensional
space spanned by $x_1y_1,x_2y_1$. The cokernel is canonically restriction to
the line $\ker P\subset G$ and is one-dimensional.

\section{Explicit cocycles and the product normal forms}
For any group with quotient coordinates $(x,y)\in\Fp^2$ put
\[
a_x(g,h,k)=x_g\Big\lfloor\frac{x_h+x_k}{p}\Big\rfloor,\quad
a_y(g,h,k)=y_g\Big\lfloor\frac{y_h+y_k}{p}\Big\rfloor,
\]
and $a_{xy}(g,h,k)=x_g\lfloor(y_h+y_k)/p\rfloor$, using representatives
$0,\ldots,p-1$. These exponents give normalized three-cocycles.
For $Q=aX^2+bXY+dY^2$ let $\alpha_Q=\zeta^{aa_x+ba_{xy}+da_y}$.
They represent all of $H^3(C_p^2,\C^\times)$; restriction to a cyclic subgroup
of generator $(u,v)$ has invariant $\zeta^{Q(u,v)}$.

In original coordinates the Heisenberg group $H_p$ has law
$(x,y,z)(u,v,w)=(x+u,y+v,z+w+xv)$. The exponent
\begin{equation}\label{eq:phi}
\phi(g,h,k)=z_gx_hy_k+\tfrac12x_g^2y_hy_k+x_gx_hy_hy_k
\end{equation}
is a three-cocycle. Indeed $\delta z=-x\smile y$, and
\[
\delta(\tfrac12x^2\smile y\smile y+x\smile(xy)\smile y)
=(x\smile y)^2.
\]
Its central slant is trivialized by $t_a(g)=\zeta^{-a z_g}$, and its central
tensor corrections vanish.

\subsection{An equivariant Heisenberg class, including $p=3$}
For $p>3$ use BCH coordinates with law
$(g,z)(h,w)=(g+h,z+w+E(g,h)/2)$, where $E(g,h)=\det(g,h)$ on $\Fp^2$.
Set $F_1=E(g,h)E(h,k)$, $F_2=E(g,h)E(g,k)$, and
\begin{equation}\label{eq:chi-large}
\chi_p(g,h,k)=\tfrac12z_gE(h,k)+\tfrac1{12}(2F_1+F_2).
\end{equation}
The Pluecker identity gives $\delta(2F_1+F_2)=3E\smile E$, proving
$\delta\chi_p=0$. It transforms exactly by $\det(M)^2$ under $M\in GL_2(\Fp)$.
Its central self-braiding is $\zeta^{-a^2}$.

For $p=3$ use original coordinates and define
\begin{equation}\label{eq:chi-three}
\chi_3=\phi+a_x+a_y.
\end{equation}
This correction is essential. For a generator $g=(x,y,z)$ the cyclic
restriction invariant of $\zeta^\phi$ has exponent $-x^2y^2$.
Adding $a_x+a_y$ makes it $1$ at every nonzero $(x,y)\in\mathbb F_3^2$,
and $0$ at zero. The leading Heisenberg spectral-sequence coefficient transforms
by $\det(M)^2$. Differences with zero leading coefficient are inflated binary
quadratic classes, detected by these cyclic restrictions; see the degree-four
calculation in \cite[Section 1.4]{MMU}. Since $\det(M)^2=1$ at $p=3$, the
class of \eqref{eq:chi-three} is $GL_2(\mathbb F_3)$-invariant.
Thus for all odd $p$ one has a chosen Heisenberg class of weight $\det^2$ and
an independently transforming inflated binary quadratic summand.

Fix $c\in\{1,\epsilon\}$ and set $v=-\epsilon/c$. Define
$\omega_{c,Q}=\zeta^{c\chi_p}\alpha_Q$ on $H_p$ and inflate it to
$\Gamma_0=H_p\times C_p$, whose last coordinate is $t$.
The centre contains the central pointed factor $\cB$ on degrees $(aZ,bT)$,
with self-braiding $\zeta^{-ca^2-vb^2}$. Here $Z$ and $T$ are the central
coordinate generators. Its half-braidings are
$\zeta^{-ca z_g-vbt_g}$, with the relevant Heisenberg coordinate used in
\eqref{eq:chi-large} or \eqref{eq:chi-three}.
Set
\[
\cC_0(c,Q)=\cB'\subset Z(\Vec_{\Gamma_0}^{\omega_{c,Q}}).
\]
The form of $\cB$ is anisotropic because $-cv=\epsilon$. Its normalized Gauss
sum is $-1$, so $\xi(\cC_0)=-1$. It factors as
\begin{equation}\label{eq:product}
\cC_0(c,Q)\simeq
\bigl(\cC(C_p,q_{-c})'\subset Z(\Vec_{H_p}^{\omega_{c,Q}})\bigr)
\boxtimes\cC(C_p,q_v).
\end{equation}
Here $q_a(u)=\zeta^{a u^2}$; the first factor has dimension $p^5$ and is
nonpointed because the centre contains $\Rep(H_p)$. The displayed central
section and every character of $\Gamma_0$ have order $p$, so the $p^3$
invertibles of $\cC_0(c,Q)$ form $C_p^3$.

\begin{proposition}\label{prop:product-orbits}
The categories with zero power map have precisely the normal forms
$\cC_0(c,Q)$ in Theorem~\ref{thm:main}; their parameters in that list are unique.
\end{proposition}
\begin{proof}
The commutator norm has two square classes, normalized to $c=1,\epsilon$.
The orthogonal complement is normalized to coefficient $v=-\epsilon/c$.
With these fixed, the metric automorphisms and changes of basis on $G$ preserving
the commutator allow exactly $\det(g)=\pm1$. Proposition~\ref{prop:h1} leaves
the full binary quadratic associator parameter. The chosen $\chi_p$ makes the
remaining action ordinary congruence, with no affine shift.
Rank zero has one orbit, rank one has its two coefficient square classes,
and rank two has its nonzero determinant as invariant. Any nondegenerate
binary form represents one over $\Fp$, hence is congruent to $X^2+dY^2$ with
$d$ its determinant. Any such congruence has determinant $\pm1$. These give
exactly $p+2$ forms for each $c$.
\end{proof}

\section{The rank-one-power normal forms}
Define $\Gamma_1$ on $(x,y,z,t)\in\Fp^4$ by
\[
(x,y,z,t)(u,w,a,b)=
((x+u)\bmod p,(y+w)\bmod p,z+a+xw,t+b+\lfloor(x+u)/p\rfloor).
\]
Let $X,Y,Z,T$ be the four coordinate generators. Then $[X,Y]=Z$,
$X^p=T$, and $Z,T$ are central of order $p$.
The map $r(g)=x_g+pt_g\in\mathbb Z/p^2$ is a homomorphism.
For $c\in\{1,\epsilon\}$, $v=-\epsilon/c$ and $s\in\Fp$ set
\begin{equation}\label{eq:prime-omega}
\omega_{c,s}(g,h,k)=\zeta^{c\phi(g,h,k)
+v x_g\lfloor(r(h)+r(k))/p^2\rfloor+s a_y(g,h,k)}.
\end{equation}
The three factors are respectively inflated from $H_p$, from a cyclic
$C_{p^2}$ cocycle with coefficient $pv$, and from $C_p$. Thus this is a
normalized three-cocycle for every odd prime, including $3$.
For $(a,b)\in\Fp^2$, the central degree $(0,0,a,b)$ has half-braiding
\begin{equation}\label{eq:tau-prime}
\tau_{a,b}(g)=\zeta^{-ca z_g}\eta^{-vb r(g)}.
\end{equation}
Its slant identity follows from $\delta z=-x\smile y$ and the homomorphism $r$.
The central tensor correction is
$\zeta^{v x_g\lfloor(b+b')/p\rfloor}$, which makes these central objects
form an order-$p^2$ elementary abelian group. Their self-braiding is
$\zeta^{-ca^2-vb^2}$. Let $\cB$ denote this anisotropic modular factor and set
\[
\cC_1(c,s)=\cB'\subset Z(\Vec_{\Gamma_1}^{\omega_{c,s}}).
\]
It has dimension $p^6$, central charge $-1$, and is nonpointed: otherwise its
product with $\cB$ would be pointed, whereas the centre contains
nonpointed $\Rep(\Gamma_1)$.

\begin{proposition}\label{prop:prime-orbits}
All rank-one-power categories have exactly the normal forms $\cC_1(c,s)$,
with distinct pairs $(c,s)$ inequivalent.
\end{proposition}
\begin{proof}[Proof retaining affine identifications]
Orthogonality puts the power image on the line perpendicular to the
commutator. Changes of basis on $G$ and metric basis on $A$ normalize the data
to $\kappa=Z e+T y_1$, $Q_A(Z)=c$, $Q_A(T)=v=-\epsilon/c$.
There are two possibilities for the commutator norm square class.
By Proposition~\ref{prop:h1}, the associator torsor modulo lower-data
automorphisms is restriction to $K=\ker P=\langle Y\rangle$, a copy of $\Fp$.
The explicit base \eqref{eq:prime-omega} with $s=0$ is identically trivial
on the entire preimage $\widetilde K=\{x=0\}\subset\Gamma_1$, since every term
of $\phi$ and of the cyclic $p^2$ factor contains an $x$ coordinate. Its
restriction to the section $\langle Y\rangle$ therefore identifies $s$ with
the residual cyclic restriction class.

The stabilizer of the normalized lower data sends
$X\mapsto aX+bY$, $Y\mapsto dY$ on $G$. Preservation of the norms of
$P(X)=T$ and $[X,Y]=Z$ forces $a^2=d^2=1$.
It therefore acts trivially on $H^3(K,\C^\times)$, since cyclic restriction
scales by $d^2$. To check the affine part as well, any crossed equivalence
induces an isomorphism of the pointed degree groups preserving $A$, the
commutator and the power map, hence preserving $\widetilde K$. Pullback of
the base cocycle remains identically trivial on $\widetilde K$. Restrict its
cochain comparison further to the section $\langle Y\rangle$: inflation from
$K$ is detected there, and the residual parameters must satisfy
$s'=d^2s=s$. Changing the lift of $Y$ does not alter this argument, since the
base vanishes on all of $\widetilde K$. The $H^1$ action also restricts to zero.
Thus no affine translation remains. Inflation of $\zeta^{s a_y}$ realizes
every $s$, proving both exhaustion and distinctness.
\end{proof}

\begin{proposition}\label{prop:prime}
$\Inv\cC_1(c,s)=C_{p^2}\times C_p$, all order-$p$ invertibles have twist one,
and $\cC_1(c,s)$ is prime.
\end{proposition}
\begin{proof}
All invertibles of the full centre have central flux $(a,b)$ and half-braiding
\eqref{eq:tau-prime} times a character $\chi$, where
$\chi(X)=\eta^k$, $\chi(Y)=\zeta^\ell$,
$k\in\mathbb Z/p^2$, $\ell\in\Fp$.
Monodromy with all $(u,w)$ of the embedded factor is trivial precisely when
$a=0$ and $k\equiv2vb\pmod p$.
The resulting group has order $p^3$. For $b\ne0$ its $p$th power is the
nontrivial character with value $\zeta^k$ on $X$, since the section
\eqref{eq:tau-prime} itself has order $p$. It is therefore $C_{p^2}\times C_p$.
Its twists are $\zeta^{vb^2}$, so every order-$p$ invertible has $b=0$ and
twist one.
A nontrivial modular factorization must have a pointed factor: two nonpointed
factors would each have dimension at least $p^5$ by \cite{CPS}. That pointed
factor can have dimension only $p$, again by the same lower bound on the other
factor. Such a factor requires an order-$p$ invertible with nontrivial twist,
which does not exist. Hence $\cC_1(c,s)$ is prime.
\end{proof}

\begin{corollary}\label{cor:exponent}
The product models have Frobenius--Schur exponent $p^2$ and the prime models
have Frobenius--Schur exponent $p^3$.
\end{corollary}
\begin{proof}
The product assertion follows from \cite[Proposition 2]{CPS} for the
dimension-$p^5$ factor and the exponent $p$ of its pointed factor.
For a pointed category $\Vec_\Gamma^\omega$, the exponent of its centre is
\[
\mathop{\rm lcm}_{\langle g\rangle\leq\Gamma}
|g|\,\operatorname{ord}(\omega|_{\langle g\rangle}),
\]
by \cite[Theorem 9.2]{NS}. In $\Gamma_1$, an element has order $p^2$
exactly when its $x$ coordinate is nonzero. On such a cyclic subgroup,
the Heisenberg and $y$ factors in \eqref{eq:prime-omega} inflate from $C_p$
and become cohomologically trivial: inflation on $H^4(-,\mathbb Z)$
multiplies the generator by $p^2$. The remaining cyclic factor has coefficient
$pv\,r(g)^2$ modulo $p^2$ and therefore order $p$. An order-$p$ cyclic
subgroup contributes at most $p^2$. Thus the full centre has exponent $p^3$.
The stabilizing factor has exponent $p$, so its centralizer has exponent $p^3$.
\end{proof}

The two branches are inequivalent by their invertible groups. The first factor
in \eqref{eq:product} is prime for the same minimum-dimension reason.
Propositions~\ref{prop:product-orbits} and \ref{prop:prime-orbits}, together
with the canonical-subcategory lemma and obstruction calculation, prove
Theorem~\ref{thm:main}.

\section{Completing the dimension-$p^6$ classification with prior branches}
For charge one, DGNO identifies every category with $Z(\Vec_K^\omega)$,
$|K|=p^3$. The complete pointed Morita classification of \cite{MMU} gives
the known $5p+32$ classes. These are prior and are not counted as new here.
Pointed categories of charge $-1$ have the classical odd-primary metric-group
normal forms: choose integers $n_r\ge0$ with $\sum r n_r=6$, and for each
$n_r>0$ a determinant square class $\delta_r\in\{1,\epsilon\}$; on
$(\mathbb Z/p^r)^{n_r}$ use diagonal coefficients
$1,\ldots,1,\delta_r$. Take the orthogonal product, retaining precisely
those with Gauss sum
\[
\prod_{r\text{ odd}}\varepsilon_p^{n_r}
\left(\frac{\delta_r}{p}\right)=-1,
\qquad \varepsilon_p=\begin{cases}1&p\equiv1\pmod4,\\i&p\equiv3\pmod4.\end{cases}
\]
Odd-primary diagonalization and its determinant criterion follow by ordinary
Gram--Schmidt over $\mathbb Z/p^r$, splitting a summand of maximum exponent
and inducting. This branch is classical, not a new result.

\section{Reproducible checks}
The supplementary repository
\href{https://github.com/numina-agents/odd-prime-dimension-six}{\texttt{odd-prime-dimension-six}} contains the
explicit check scripts and exact source provenance. The scripts
\texttt{check\_control.py} and \texttt{check\_prime\_control.py} check all
$27^4$ and $81^4$ cocycle equations at $p=3$, together with central slants,
tensor corrections and self-braidings. The script \texttt{check\_chi.py}
checks the exceptional cyclic restrictions and the Pluecker primitive.
A separate symbolic check verifies the polynomial cocycle and central-slant
identities over $\mathbb Q$. These are consistency checks of the explicit
cochains; the equivalence and exhaustion proofs do not rely on a census.

\subsection*{Acknowledgments}
We thank Phoebe for independent checks of the central-flux and canonical
Tannakian reductions, and Cygnus for the complete independent assessment,
the symbolic verification and the clearer preimage-$K$ exposition of the
affine restriction argument. These results have not been checked by a human.

\begin{thebibliography}{99}
\bibitem{ENO05} P. Etingof, D. Nikshych, and V. Ostrik,
\emph{On fusion categories}, Ann. of Math. 162 (2005), 581--642,
arXiv:math/0203060.
\bibitem{GN} J. Green and D. Nikshych,
\emph{The Tannakian radical and the mantle of a braided fusion category},
arXiv:2412.16722v2 (2025).
\bibitem{Schopieray} A. Schopieray,
\emph{Nondegenerate extensions of near-group braided fusion categories},
Rev. Uni\'on Mat. Argentina 64 (2023), 413--438, arXiv:2109.04011.
\bibitem{Naidu} D. Naidu, \emph{Crossed pointed categories and their
equivariantizations}, Pacific J. Math. 247 (2010), 477--496.
\bibitem{MN} G. Mason and S.-H. Ng,
\emph{Cleft extensions and quotients of twisted quantum doubles},
Developments in Mathematics 38 (2014), 229--246, arXiv:1404.2016.
\bibitem{NN} D. Naidu and D. Nikshych,
\emph{Lagrangian subcategories and braided tensor equivalences of twisted
quantum doubles of finite groups}, arXiv:0705.0665.
\bibitem{CPS} A. Czenky, J. Plavnik, and A. Schopieray,
\emph{On Frobenius--Schur exponent bounds}, arXiv:2404.06643, Section 4.2.
\bibitem{DGNO} V. Drinfeld, S. Gelaki, D. Nikshych, and V. Ostrik,
\emph{Group-theoretical properties of nilpotent modular categories},
arXiv:0704.0195.
\bibitem{GR} T. Gannon and A. Riesen, \emph{Orbifolds of Pointed Vertex Operator
Algebras I}, arXiv:2410.00809.
\bibitem{MMU} K. Maya, A. Mej\'ia Casta\~no, and B. Uribe,
\emph{Classification of pointed fusion categories of dimension $p^3$ up to weak
Morita equivalence}, arXiv:1808.05139.
\bibitem{ENO} P. Etingof, D. Nikshych, and V. Ostrik,
\emph{Fusion categories and homotopy theory}, arXiv:0909.3140.
\bibitem{DN} A. Davydov and D. Nikshych,
\emph{Braided Picard groups and graded extensions of braided tensor categories},
arXiv:2006.08022.
\bibitem{NS} S.-H. Ng and P. Schauenburg,
\emph{Frobenius--Schur indicators and exponents of spherical categories},
arXiv:math/0601012, Theorem 9.2.
\end{thebibliography}
\end{document}
